% Working proof notes for the parent report, not a standalone final artifact.
% Exact comparison point: ProveIt 8662bc0a00f4e77e482f7d141d7a4f919f70f9aa.
% Environments needed: theorem, lemma, proposition, corollary.

\section{Putting induction at the closure it actually proves}

The change in this section is a refactoring of the primitive basis, not a
new induction principle.  The baseline already proves the new rule as
\texttt{NPrf.qIProg} in \texttt{NativeEq.lean}.  The point is to take that
rule as primitive and delete the existentially wrapped progress rule.
This removes a witness and one substitution from the primitive induction
interface.  The statement below concerns the mathematical calculus and
does not depend on an executable checker's correctness.

\subsection{Two complete calculi in exactly the same language}

Both calculi use numerical terms
\[
 t ::= x\mid 0\mid S t\mid I(n,a;u.h),
\]
where only the occurrence of $u$ in the fixed finite body $h$ is bound.
Native formulas have grammar
\[
 A ::= (s=t)\mid Q_y A\mid D_{x,y}A.
\]
The two displayed bound variables of $D$ are distinct.  All substitutions
avoid capture; bound variables may be renamed.  Contexts are finite
reusable lists of arbitrary native formulas.  In particular, they are
not restricted to equations or closed formulas.

The common computation axioms, in every context, are
\[
 I(0,a;u.h)=a,\qquad
 I(Sn,a;u.h)=h[I(n,a;u.h)/u].
\]
The following inference rules are also common:
\[
 \frac{}{\Gamma\vdash A}\ \ (A\in\Gamma),
 \qquad
 \frac{\Gamma\vdash s=t\qquad\Gamma\vdash F(s)}
      {\Gamma\vdash F(t)}\ \mathrm{Tr},
\]
\[
 \frac{\Gamma\vdash Q_yB(y)\qquad\Gamma\vdash B(0)}
      {\Gamma\vdash B(r)}\ \mathrm{QE},
\]
\[
 \frac{\Gamma\vdash D_{x,y}B(x,y)\qquad
       \Gamma,Q_yB(b,y)\vdash A}
      {\Gamma\vdash A}\ \mathrm{DE},
\]
\[
 \frac{\Gamma,B(t,0)\vdash B(t,Sz)}
      {\Gamma\vdash D_{x,y}B(x,y)}\ \mathrm{DI}_{\rm case}.
\]
In Tr, $F$ is any fixed native formula template, with an arbitrary finite
number of occurrences of its numerical placeholder.  In DE, $b$ is fresh
for $\Gamma,A$ and all passive parameters of $B$.  In DI$_{\rm case}$,
$z$ is fresh for $\Gamma,t$ and all passive parameters of $B$.  The
displayed local assumptions are discharged.  There are no additional
logical or arithmetic axioms, and no primitive reflexivity rule.

The old calculus $\mathcal C_D$ additionally has
\[
 \frac{\Gamma,E(t,z)\vdash E(t,Sz)}
      {\Gamma\vdash D_{x,y}E(x,y)}\ \mathrm{DI}_{\rm prog},
 \tag{DP}
\]
where $E(x,y)$ is a literal equation between numerical terms and $z$ has
the same freshness condition as in DI$_{\rm case}$.

The new calculus $\mathcal C_Q$ replaces (DP) by
\[
 \boxed{\quad
 \frac{\Gamma,E(z)\vdash E(Sz)}
      {\Gamma\vdash Q_yE(y)}\ \mathrm{QI}_{\rm prog}.
 \quad}
 \tag{QP}
\]
Here $E(y)$ is a literal equation between numerical terms, and $z$ is
fresh for $\Gamma$ and every passive parameter of $E$.  The arbitrary
parameters of $E$ need not be closed.  The restriction concerns the root
of its finite syntax: a $Q$ or $D$ formula containing an equation is not
a literal equation.  There is no witness field in (QP).

\subsection{Structural facts used by the proof}

\begin{lemma}
In either calculus, weakening, exchange, contraction, injective numerical
renaming, capture-free numerical substitution, and ordinary-assumption
cut are admissible.
\end{lemma}
\begin{proof}
Induct on a finite derivation.  Before inserting an additional context,
term, or proof, rename the finitely many local eigenparameters away from
its free variables.  Every common rule is preserved.  For (DP) and (QP),
numerical substitution preserves the literal-equation root, and the fresh
step parameter can be kept fixed after that renaming.  For assumption
cut, replace each use of the cut hypothesis by the supplied derivation,
weakened to its current local context.  Eigenparameters are freshened
before this replacement.  No arithmetic induction or semantic argument
is used in these external inductions on derivations.
\end{proof}

Only weakening and renaming are needed for the equivalence construction
below.  In particular, equality transport is not being used to change
the rule being simulated.

\subsection{Identity-sequent equivalence}

\begin{theorem}[Direct closure progress]
For every finite context $\Gamma$ of arbitrary native formulas and every
native formula $A$,
\[
       \Gamma\vdash_{\mathcal C_D} A
       \quad\Longleftrightarrow\quad
       \Gamma\vdash_{\mathcal C_Q} A.
\]
Both implications are effective transformations of finite derivations.
They leave the terms, formulas, context, and final conclusion unchanged.
\end{theorem}
\begin{proof}
We first simulate (QP) in $\mathcal C_D$.  Suppose its transformed premise
is
\[
                \pi:\quad\Gamma,E(z)\vdash E(Sz).
\]
Choose a variable $u$ absent from $E$ and $\Gamma$.  Regard
$\widetilde E(u,y):=E(y)$ as an equation with an unused witness coordinate.
Apply (DP) with witness $0$ and premise $\pi$ to obtain
\[
                    \Gamma\vdash D_{u,y}E(y).
\]
Use DE.  Its fresh witness parameter $b$ does not occur in the opened
closure, which is exactly $Q_yE(y)$.  Its continuation is therefore the
hypothesis derivation
\[
              \Gamma,Q_yE(y)\vdash Q_yE(y).
\]
The DE conclusion is the required $\Gamma\vdash Q_yE(y)$.
Fresh $u,b$ exist because all syntax and contexts are finite.

Conversely, suppose an old (DP) instance has recursively transformed
premise
\[
                \pi:\quad\Gamma,E(t,z)\vdash E(t,Sz).
\]
Set $F(y):=E(t,y)$, first renaming bound variables to avoid capture.
Since $z$ is absent from $t$ and all retained data, $\pi$ is a legitimate
(QP) premise for this fixed equation template.  It yields
\[
                         \Gamma\vdash Q_yE(t,y).
\]
Now use DI$_{\rm case}$ with the original witness $t$.  Let $w$ be fresh
for $\Gamma,t$ and the passive parameters of $E$, and work under its
local assumption $E(t,0)$.  Weaken the displayed closure proof into this
context.  QE, applied at $Sw$ and supplied with the base hypothesis,
gives
\[
              \Gamma,E(t,0)\vdash E(t,Sw).
\]
DI$_{\rm case}$ discharges that assumption and gives exactly the old
conclusion $\Gamma\vdash D_{x,y}E(x,y)$.  If $w$ occurs as an internal
eigenparameter in $\pi$, rename that internal parameter before weakening.
No equality theorem, computation equation, or numerical substitution into
the proof $\pi$ is needed: $F$ is the equation already appearing in its
premise and conclusion.

Induction on the entire old or new derivation now replaces the one
different primitive rule by its displayed simulation.  Common rules are
copied after recursively translating their premises.  This establishes
both implications for arbitrary contexts and arbitrary native goals.
\end{proof}

\begin{proposition}[Inference-node overhead]
Count every hypothesis, computation axiom, and primitive rule occurrence
as one node.  An $N$-node new proof with $r$ occurrences of (QP) translates
to an old proof with $N+2r\leq3N$ nodes.  An $M$-node old proof with $s$
occurrences of (DP) translates to a new proof with $M+3s\leq4M$ nodes.
\end{proposition}
\begin{proof}
The first replacement has DP, DE, and Hyp, and inserts the old premise
once.  The second replacement has QP, QE, Hyp, and DI$_{\rm case}$, and
also inserts its original premise once.  Weakening and fresh renaming
do not change inference-node counts.
\end{proof}
This is a proof-tree bound, not a bound on serialized formula sizes,
certificate bytes, checker runtime, or program-source size.

\subsection{Induction and the ordinary interpretation}

The ordinary reading remains
\[
 (Q_yB(y))^\circ=B^\circ(0)\to\forall y\,B^\circ(y),
 \qquad
 (D_{x,y}B)^\circ=
     \exists x\bigl(B^\circ(x,0)\to\forall y\,B^\circ(x,y)\bigr).
\]
Numerical terms and equality retain their meanings.  From a proof of
$E(0)$ and a proof of $E(Sz)$ under $E(z)$, (QP) followed by QE gives
$E(r)$ for any numerical term $r$.  Thus the new arithmetic interface
is contextual equation induction in its conditional-base form.

For direct soundness in ordinary $T_{\rm eq}$, assume the translation of
the retained context and then assume $E(0)$.  The translated premise
of (QP) supplies the step under $E(z)$; contextual equation induction
gives $\forall y E(y)$.  Discharge the base assumption.  The result is
precisely the stated ordinary reading of $Q_yE(y)$.  The step
eigenparameter is absent from the context and from $E(0)$, as required.

If (QP) were allowed for arbitrary matrices $B$, the same two-rule
argument with QE would give induction for every native formula.  Since
the native logical macros represent all first-order formulas, this would
permit unrestricted first-order induction.  The literal-equation guard
must therefore remain explicit; moving induction from $D$ to $Q$ does
not license weakening that guard.

\begin{corollary}[Transfer of the full iPRA bi-interpretation]
The supplied ordinary/native interpretation and its inverse are also
interpretations for $\mathcal C_Q$, with the same formula translations
and the same proved round trips.  Composing them with the scalar
$T_{\rm eq}$--iPRA bi-interpretation gives an equiconsistent
bi-interpretation between $\mathcal C_Q$ and full first-order iPRA.
Both numerical domains remain whole, equality is actual equality, and
the numerical return isomorphisms retain graph $x=y$.
\end{corollary}
\begin{proof}
The old and new native systems have identical finite syntax and, by the
theorem, identical consequence relations in arbitrary contexts.  Any
derivation verifying an axiom translation, inference translation, or
formula round trip in the old native system therefore transfers to the
new one.  The identity maps between the two native systems compose to
the literal identity on syntax.  They do not change the scalar-to-iPRA
term compiler, its inverse, or the internal equalities proving their
composites.  Transfer applies to arbitrary formulas, not merely closed
equations or equiconsistency assertions.
\end{proof}
The argument uses the written scalar/iPRA bridge, which is part of the
baseline mathematical development.  The source also contains
\texttt{BiInterpretation.lean} for the scalar/ordinary iPRA bridge.
These observations are provenance, not a claim that a new executable
checker or this entire new basis has already been machine verified.

\subsection{What this refactoring does not make redundant}

The equivalence proof is independent of Tr and of both iterator
computation axioms: it works over the Hyp/QE/DE/DI$_{\rm case}$ skeleton,
with any common collection of additional rules preserving the structural
operations used.  Accordingly, direct Q-progress does not by itself
make a previously inadequate restriction of equality transport adequate.
Any such further simplification needs a separate argument.  In
particular, it does not evade the supplied countermodels for deleting
essential pieces of equality replacement.

Nor can DI$_{\rm case}$ simply be removed from the new calculus.  In the
remaining system, starting with a context containing only equations and
$Q$ formulas whose matrices are equations, every derivable formula is
again of one of those two forms.  Prove this by induction on a derivation:
Hyp and computation have that property, QP introduces precisely a $Q$
of an equation, QE returns an equation, and numerical transport preserves
the formula-constructor shape.  A DE node is impossible because its
major would be a derived $D$ formula.  Hence, in particular, no $D$
formula is a theorem in that case-free remainder.  By contrast,
$D_{x,y}(0=0)$ is proved in the full system by DI$_{\rm case}$ at witness
$0$, using its local hypothesis $0=0$ as its successor premise.

The modest simplification is therefore precisely located: the induction
primitive proves a closure directly, without an irrelevant existential
witness.  Neither a new logical connective, a new numerical operation,
nor a stronger induction scheme has been introduced.  Any claim of
smaller complete checker size must be supported by the complete program
and its measured, validated public input contract.

\section{Optional deterministic introduction at every closure goal}

For a particularly regular implementation, extend the preceding new
calculus by a directly derived case introduction for compound $Q$ goals.
Equivalently, replace (QP) by the single rule schema
\[
 \frac{\Gamma,B(\eta_B(z))\vdash B(Sz)}
      {\Gamma\vdash Q_y B(y)}\ \mathrm{QI},
 \qquad
 \eta_B(z)=
 \begin{cases}
 z,& B\text{ has a literal-equation root},\\
 0,& B\text{ has a }Q\text{ or }D\text{ root}.
 \end{cases}
 \tag{DQI}
\]
The same freshness condition applies: $z$ is absent from the retained
context and all passive parameters of $B$.  Keep DI$_{\rm case}$ as the
only primitive $D$ introduction.  Call the resulting calculus
$\mathcal C_{\rm det}$.  A proof certificate supplies no independent
case/progress selector: the requested outer constructor and, for a $Q$
goal, its matrix root determine which local assumption is used.

\begin{theorem}[Deterministic goal-directed introduction]
For every context $\Gamma$ and native formula $A$,
\[
 \Gamma\vdash_{\mathcal C_D}A
 \quad\Longleftrightarrow\quad
 \Gamma\vdash_{\mathcal C_Q}A
 \quad\Longleftrightarrow\quad
 \Gamma\vdash_{\mathcal C_{\rm det}}A.
\]
All three systems have identical finite syntax and final sequents.
The inference-node bounds $3N$ from the deterministic calculus to the old
calculus and $4M$ in the converse direction remain valid.
\end{theorem}
\begin{proof}
An equation instance of (DQI) is (QP), already simulated.  For a compound
matrix, its premise is $\Gamma,B(0)\vdash B(Sz)$.  Choose an unused
witness coordinate $u$.  DI$_{\rm case}$ with witness $0$ yields
$D_{u,y}B(y)$, and DE with the hypothesis continuation returns $Q_yB(y)$.
This is the same three-node construction, using D-case rather than
D-progress.  It is valid for any matrix, including a literal equation.
Thus every newly added compound-Q introduction is already derivable in
both prior calculi.  The previous old-to-new simulation uses only the
equation branch of (DQI), so it is unchanged.
\end{proof}

The last observation explains why this change is not the invalid proposal
that uses only progress introduction whenever a $D$ matrix is an equation.
In the deterministic calculus, a case introduction of $Q_yE(y)$ for an
equation is still derivable through the dummy-witness $D$ construction.
For example, $Q_y(y=1)$ has a false zero instance and is a theorem of the
intended native calculus, even though its direct Q-progress premise
$(z=1)\vdash(Sz=1)$ is invalid.  Its legitimate zero-case construction
uses D-case; the distinction has not been erased.

Removing D-case still loses every closed theorem with a $D$ root, even
in $\mathcal C_{\rm det}$.  More generally, if a context contains no
formula with a $D$ root, no $D$-root formula is derivable after that
deletion.  Induct on a putative derivation of such a conclusion.
Hyp is unavailable; computation and QI have other roots; Tr requires
a minor with the same $D$ root; QE requires a base minor with the same
$D$ root; and DE requires a major with a $D$ root.  Each applicable
case contradicts the induction hypothesis on the indicated shorter
premise.  Assumptions under a QI subderivation do not invalidate this
argument, because QI itself cannot be the last rule of a $D$-root goal.

The deterministic extension preserves the ordinary interpretation and
the full iPRA bi-interpretation exactly as above.  Its equation branch
is sound by equation induction, and its compound branch is sound by
numerical coverage followed by implication introduction.  Its interest
is a simpler primitive proof interface with no independent introduction
mode, not a new semantic principle.  In particular, its concrete size
advantage, if any, must be measured on a complete sound checker.
